% ============================================================
\chapter{Continuous-Time Markov Chains}
\label{ch:continuous_markov}
% ============================================================

The passage from discrete to continuous time is not a mere formality: it profoundly transforms the nature of Markov chains. In discrete time, transitions occur at regular instants. In continuous time, the process may jump from one state to another at any moment, and the sojourn time in each state follows an exponential distribution---the only continuous memoryless distribution. This link between the Markov property and the exponential law is one of the most elegant observations in the theory. Continuous-time Markov chains model queues (M/M/1 systems), stochastic chemical reactions, telecommunications networks, and many other systems where events occur randomly in time.

\section{Definition and the Markov Property}

\begin{definition}[Continuous-Time Markov Chain]
Let $E$ be a countable state space. A process $(X_t)_{t \geq 0}$ with values
in $E$ and c\`adl\`ag sample paths is a \emph{continuous-time Markov chain} (CTMC)
if for all $s, t \geq 0$ and every $j \in E$:
\[
  \PP(X_{t+s} = j \mid \mathcal{F}_s) = \PP(X_{t+s} = j \mid X_s) \quad \text{a.s.}
\]
The chain is \emph{homogeneous} if $\PP(X_{t+s} = j \mid X_s = i) = p_{ij}(t)$
does not depend on $s$.
\end{definition}

\begin{definition}[Transition Semigroup]
The family of matrices $(P(t))_{t \geq 0}$ with $P(t) = (p_{ij}(t))_{i,j \in E}$
is called the \emph{transition semigroup}. It satisfies:
\begin{enumerate}
  \item $P(0) = I$ (identity matrix),
  \item $P(s+t) = P(s)P(t)$ for all $s, t \geq 0$ (semigroup property),
  \item $p_{ij}(t) \geq 0$ and $\sum_j p_{ij}(t) = 1$ for all $i$.
\end{enumerate}
\end{definition}

\section{Infinitesimal Generator}

\begin{definition}[Infinitesimal Generator]
The \emph{infinitesimal generator} (or \emph{rate matrix}) of the CTMC is the
matrix $Q = (q_{ij})_{i,j \in E}$ defined by:
\[
  Q = \lim_{t \to 0^+} \frac{P(t) - I}{t}, \quad \text{i.e.} \quad
  q_{ij} = \lim_{t \to 0^+} \frac{p_{ij}(t) - \delta_{ij}}{t}.
\]
\end{definition}

\begin{proposition}[Properties of $Q$]
The matrix $Q$ satisfies:
\begin{enumerate}
  \item $q_{ij} \geq 0$ for $i \neq j$,
  \item $q_{ii} \leq 0$ for all $i$,
  \item $\sum_{j \in E} q_{ij} = 0$ for all $i$, i.e.\ $q_{ii} = -\sum_{j \neq i} q_{ij}$.
\end{enumerate}
We write $q_i = -q_{ii} = \sum_{j \neq i} q_{ij}$ for the \emph{exit rate} from
state $i$. If $q_i > 0$, the holding time in $i$ is exponentially distributed
with parameter $q_i$.
\end{proposition}

\begin{intuition}[title=\textbf{Construction via Holding Times}]
A CTMC can be constructed as follows:
\begin{enumerate}
  \item Starting from state $i$, the process remains there for an exponential
    time with parameter $q_i$.
  \item At the end of this time, it jumps to $j \neq i$ with probability
    $q_{ij}/q_i$.
  \item The process restarts from the new state.
\end{enumerate}
The embedded discrete chain (the sequence of visited states) has transition
matrix $\tilde{P}$ with $\tilde{p}_{ij} = q_{ij}/q_i$ for $j \neq i$.
\end{intuition}

\section{Kolmogorov Equations}

\begin{theorem}[Kolmogorov Forward Equation]
\label{thm:kolmo_forward}
Under regularity conditions (e.g.\ $\sup_i q_i < \infty$):
\[
  P'(t) = P(t) Q, \quad \text{i.e.} \quad
  \frac{d}{dt} p_{ij}(t) = \sum_{k \in E} p_{ik}(t) q_{kj}.
\]
\end{theorem}

\begin{theorem}[Kolmogorov Backward Equation]
\label{thm:kolmo_backward}
Under the same conditions:
\[
  P'(t) = Q P(t), \quad \text{i.e.} \quad
  \frac{d}{dt} p_{ij}(t) = \sum_{k \in E} q_{ik} p_{kj}(t).
\]
\end{theorem}

\begin{remark}
If $E$ is finite, the solution is $P(t) = e^{Qt}$, where the matrix exponential
is defined by $e^{Qt} = \sum_{n=0}^{\infty} \frac{(Qt)^n}{n!}$.
\end{remark}

\begin{keyformulas}[title=\textbf{Kolmogorov Equations}]
\begin{align*}
  &\text{Backward:} & P'(t) &= QP(t), \quad P(0) = I, \\
  &\text{Forward:} & P'(t) &= P(t)Q, \quad P(0) = I.
\end{align*}
Solution (finite state space): $P(t) = e^{Qt}$.
\end{keyformulas}

\section{Stationary Distributions and Convergence}

\begin{definition}[Stationary Distribution]
A probability vector $\pi = (\pi_i)_{i \in E}$ is a \emph{stationary distribution}
for the CTMC if:
\[
  \pi Q = 0, \quad \text{i.e.} \quad \sum_{i \in E} \pi_i q_{ij} = 0
  \quad \forall j \in E.
\]
Equivalently, $\pi P(t) = \pi$ for all $t \geq 0$.
\end{definition}

\begin{theorem}[Convergence to Equilibrium]
\label{thm:convergence_equilibrium}
If the CTMC is irreducible and positive recurrent, there exists a unique
stationary distribution $\pi$, and for all $i, j \in E$:
\[
  \lim_{t \to \infty} p_{ij}(t) = \pi_j.
\]
\end{theorem}

\section{Birth-and-Death Processes}

\begin{definition}[Birth-and-Death Process]
A birth-and-death process on $E = \N$ is a CTMC with generator:
\[
  q_{i,i+1} = \lambda_i \quad (\text{birth rate}), \quad
  q_{i,i-1} = \mu_i \quad (\text{death rate}), \quad i \geq 0,
\]
with $\mu_0 = 0$, $\lambda_i > 0$ for $i \geq 0$, $\mu_i > 0$ for $i \geq 1$,
and $q_{ij} = 0$ for $\abs{i - j} > 1$.
\end{definition}

\begin{theorem}[Stationary Distribution]
\label{thm:bd_stationary}
The birth-and-death process admits a stationary distribution if and only if
\[
  S = \sum_{n=1}^{\infty} \frac{\lambda_0 \lambda_1 \cdots \lambda_{n-1}}
  {\mu_1 \mu_2 \cdots \mu_n} < \infty.
\]
In this case, $\pi_0 = (1 + S)^{-1}$ and
$\pi_n = \pi_0 \frac{\lambda_0 \cdots \lambda_{n-1}}{\mu_1 \cdots \mu_n}$.
\end{theorem}

\begin{proof}
The detailed balance equations $\pi_i q_{i,i+1} = \pi_{i+1} q_{i+1,i}$
give $\pi_i \lambda_i = \pi_{i+1} \mu_{i+1}$, whence the formula by induction.
Normalisability is equivalent to $S < \infty$.
\end{proof}

\begin{example}[$M/M/1$ Queue]
Poisson arrivals at rate $\lambda$, exponential service at rate $\mu$:
$\lambda_i = \lambda$, $\mu_i = \mu$ for all $i$. The stability condition is
$\rho = \lambda/\mu < 1$, and the stationary distribution is geometric:
$\pi_n = (1 - \rho)\rho^n$.
\end{example}

\begin{example}[$M/M/c$ Queue]
With $c$ servers: $\lambda_i = \lambda$ for all $i$ and
$\mu_i = \min(i, c) \mu$. The stability condition is $\rho = \lambda/(c\mu) < 1$.
\end{example}

\begin{example}[$M/M/\infty$ Queue]
Infinitely many servers: $\lambda_i = \lambda$, $\mu_i = i\mu$. The stationary
distribution is Poisson: $\pi_n = e^{-\lambda/\mu} \frac{(\lambda/\mu)^n}{n!}$.
\end{example}

\section{Explosion and Regularity}

\begin{definition}[Explosion]
A CTMC \emph{explodes} if the number of jumps in finite time is infinite, i.e.\
$T_{\infty} = \sum_{n=0}^{\infty} S_n < \infty$ a.s., where $S_n$ is the $n$-th
holding time.
\end{definition}

\begin{theorem}[Non-Explosion Criterion]
\label{thm:non_explosion}
If $\sup_{i \in E} q_i < \infty$ (bounded exit rates), the CTMC does not explode.
More generally, the chain does not explode if and only if the minimal solution
of the backward Kolmogorov equations is stochastic (i.e.\ $\sum_j p_{ij}(t) = 1$).
\end{theorem}

\begin{example}[Yule Process]
The pure birth process with $\lambda_n = n\lambda$ (each individual reproduces
independently at rate $\lambda$). We have $\E[X_t] = X_0 e^{\lambda t}$, and the
chain does not explode since $\sum_n \frac{1}{q_n} = \sum_n \frac{1}{n\lambda} = +\infty$.
\end{example}

\section{Strong Markov Property}

\begin{theorem}[Strong Markov Property --- Continuous Time]
\label{thm:strong_markov_ct}
Let $(X_t)_{t \geq 0}$ be a regular (non-explosive) CTMC and $\tau$ an a.s.\
finite stopping time. Then, conditionally on $X_{\tau} = i$, the process
$(X_{\tau + t})_{t \geq 0}$ is a CTMC with the same generator, starting from $i$,
independent of $\mathcal{F}_{\tau}$.
\end{theorem}

\section{Reversibility in Continuous Time}

\begin{definition}[Reversibility]
A CTMC with generator $Q$ is reversible with respect to a distribution $\pi$ if:
\[
  \pi_i q_{ij} = \pi_j q_{ji}, \quad \forall i, j \in E.
\]
\end{definition}

\begin{theorem}[Burke's Theorem]
\label{thm:burke}
In an $M/M/1$ queue in steady state ($X_0 \sim \pi$), the departure process is
a Poisson process with rate $\lambda$.
\end{theorem}

\section{Exercises}

\begin{exercise}
Consider a CTMC on $E = \{1, 2, 3\}$ with generator
\[
  Q = \begin{pmatrix} -3 & 2 & 1 \\ 1 & -2 & 1 \\ 2 & 1 & -3 \end{pmatrix}.
\]
Compute the stationary distribution. Check whether the chain is reversible.
\end{exercise}

\begin{exercise}
For a birth-and-death process with $\lambda_n = \lambda$ and $\mu_n = n\mu$
($n \geq 1$), find the stationary distribution and show that the chain does not
explode.
\end{exercise}

\begin{exercise}
Show that for a CTMC on a finite state space, the matrix exponential
$P(t) = e^{Qt}$ satisfies both the forward and backward Kolmogorov equations.
\end{exercise}

\begin{exercise}
Consider an $M/M/1$ queue with $\lambda = 2$ and $\mu = 3$. Compute the
steady-state probability that the system is empty, the average number of
customers, and the average sojourn time (Little's law).
\end{exercise}

\begin{exercise}
Show that a pure birth process with $\lambda_n = n^2$ explodes in finite time.
\emph{Hint: show that $\sum_{n} 1/\lambda_n < \infty$.}
\end{exercise}

\begin{exercise}
For the continuous-time Ehrenfest chain (each ball is moved independently at
rate $1$), write down the generator $Q$ and find the stationary distribution.
\end{exercise}
